\section{From benchmark gaps to a joint execution law}
\label{sec:realisticexecution}

The usefulness of AC and OW is that their assumptions make a decision
transparent. Their limitations should be attached to those assumptions,
not attributed to every member of an entire model family. In particular,
AC uses linear permanent impact with a constant \emph{coefficient}; the
price displacement is not constant. Under the continuous, fixed-quantity,
single-asset convention of Proposition~\ref{prop:perminvisible}, its
permanent cost is schedule-independent. The original AC paper also
considers portfolios and a liquidity-adjusted VaR criterion
\citep{almgren2001}. OW explicitly retains transient impact and book
resilience \citep{obizhaevawang2013}; its block-shaped benchmark is not a
model without liquidity recovery.

Likewise, an empirical square-root relation between aggregate metaorder
size and price change is not an instruction to replace instantaneous
trading cost by a square root. The relevant input, horizon and conditioning
must match. A concave price response can produce a convex total charge,
while a nonlinear transient specification can admit profitable round trips
unless response and decay satisfy additional restrictions
\citep{gatheral2010,curato2017}. Sections~\ref{sec:nonlinear}
and~\ref{sec:extensions} retain those distinctions.

\subsection{Six requirements and their present status}\label{sec:requirements}

\paragraph{Fill law.}
Sections~\ref{sec:types}, \ref{sec:venues}, \ref{sec:bookbridge}
and~\ref{sec:algorithm} already distinguish instructions, venues,
reservations, tagged orders and executable messages. A schedule is a
requested control, not proof that every requested share fills. A market
application still needs venue-specific priority rules, size-dependent
partial fills, tick and lot admissibility, and the joint timing of
messages, acknowledgements and executions. Replacing an order resets its
priority according to those rules; a pending cancel does not release a
reservation merely because the trader sent it.

\paragraph{Impact and changing liquidity.}
Nonlinearity and cross-impact are represented in the structural analysis;
they are not empirically calibrated here. An intraday clock, a volatility
state and a signal can enter the controlled book law. A U-shaped seasonal
profile is a candidate fit to test by venue, asset and day type, not a
universal law to impose. Own impact must enter the transition law as well
as the current transaction price if it changes subsequent opportunities.
A signal-dependent drift must also be separated from price movement
caused by the trader.

\paragraph{Information and adverse selection.}
The belief and observation certificates address discarded information and
supplied uncertainty sets. They do not make parameters known. A marked
jump can join a fill, a news shock and a deterioration of depth in one
event. Calibration must preserve their dependence. Hidden toxicity calls
for a belief state or another sufficient observation summary; enlarging a
state vector with unavailable information creates an infeasible benchmark.
A survival estimator for fill time also needs a stated censoring target:
policy-driven cancellation can be informative, and cancellation and fill
can instead be modeled as competing events. Treating cancellation as
independent censoring requires justification.

\paragraph{Competition.}
A controlled book can encode a reduced-form response to displayed size,
quote changes and schedule predictability. It does not thereby solve a
game among strategic agents. A Nash claim needs opponents' information,
action sets and best responses, and a joint equilibrium calculation.
The model-risk family should include changes in the response to one's own
policy; perturbing only exogenous order-arrival rates misses this channel.

\paragraph{Objectives and constraints.}
The mandate ledger and terminal condition are part of the problem.
Fees, rebates and hedge costs belong in net cash flows. Capital and
inventory limits belong in feasibility; drawdown needs wealth and its
running peak if modeled as a path constraint. A quadratic risk charge is
one declared preference, not a substitute for every such requirement.
Adaptive mean--variance and ratio objectives do not in general satisfy the
ordinary additive Bellman recursion. Precommitment or an equilibrium
interpretation must be specified before claiming time consistency.

\paragraph{Evaluation.}
The monograph already supplies model-error and robustness certificates. What
remains is evidence that their uncertainty family covers the relevant
market law. A fixed historical book cannot identify the market response
to an unobserved intervention merely by inserting simulated orders.
Counterfactual support, dependence stress and impact-aware evaluation are
separate from solving the chosen control problem accurately.

\subsection{What is shared by the Bellman step}

At decision time \(t\), let \(x\) retain remaining quantity, live-order
reservations, available queue information, impact memory and any observed
signal or hidden-state belief. A joint kernel
\(P_t(\dd z\mid x,a)\) describes executed sizes and prices, fees,
price/depth changes and message acknowledgements. With a state update
\(F_t\) and a one-step shortfall charge \(\ell_t\), a finite-horizon
expected-cost problem reads
\begin{equation}\label{eq:jointexecutionbellman}
 V_t(x)=\inf_{a\in\mathcal A_t(x)}
  \int\{\ell_t(x,a,z)+V_{t+1}(F_t(x,a,z))\}
                  P_t(\dd z\mid x,a).
\end{equation}
Every allowed event must preserve the physical ledger. The terminal value
specifies completion, an executable fallback, or an explicit failure
penalty. Assigning \(+\infty\) to noncompletion is a hard constraint, not a
claim that completion is always feasible.

This recursion is shared by many models because it is conditional
optimization. It does not justify independent fill and price kernels.
Nor does inserting a worst-case kernel separately at each node necessarily
evaluate one unknown, fixed model: a nodewise adversary permits different
choices after different histories. Such rectangular uncertainty is a
stronger problem unless the family is stable under this pasting, or the
state retains the common latent parameter. The baseline-relative
certificate in Section~\ref{sec:baselineexecution} instead compares the
proposal and baseline under the same model.

\subsection{A dependence term that reverses the routing decision}

A two-date completion problem isolates one missing interaction without
claiming to describe a calibrated exchange. There are \(Q\) shares left.
An immediate market action costs \(C_M\). Alternatively, a passive order
fills \(F\in[0,Q]\), earns a rebate \(r\) per share, and the residual is
completed at the next date at an incremental cost \(K\) per share. Both
choices complete the mandate; the fallback is assumed available for the
full residual. All prices are relative to one initial benchmark, with
\(K\) including whatever terminal slippage this supplied law specifies.

\begin{proposition}[The fill--completion dependence charge]
\label{prop:completiondependence}
For square-integrable \(F,K\), the expected passive cost is
\begin{equation}\label{eq:completioncovariance}
 C_P=\E K\,(Q-\E F)-r\E F-\Cov(K,F).
\end{equation}
An independent coupling with the same marginals underestimates this cost
by \(-\Cov(K,F)\) when the covariance is negative. Its absolute error is
at most \(\sqrt{\Var(K)\Var(F)}\). If only discrete marginals
\(p_i=\Pr(F=f_i)\), \(q_j=\Pr(K=k_j)\) are supplied, the exact smallest
and largest expected costs are the minimum and maximum of
\begin{equation}\label{eq:completiontransport}
 \sum_{ij}\pi_{ij}\{k_j(Q-f_i)-rf_i\},\qquad
 \pi_{ij}\ge0,\quad \sum_j\pi_{ij}=p_i,\quad \sum_i\pi_{ij}=q_j.
\end{equation}
\end{proposition}
\begin{proof}
Expand \(\E[K(Q-F)-rF]\) and use
\(\E[KF]=\E K\E F+\Cov(K,F)\). Cauchy--Schwarz gives the bound.
Every joint law with the stated finite marginals is a feasible matrix
\(\pi\), and every such matrix defines a joint law. The feasible set is
nonempty and compact, so the linear objective attains both extrema.
\end{proof}

Take \(Q=1\), \(r=0.001\), \(C_M=0.06\). A paired three-event law has
\[
 (F,K)=(1,0.02),\ (0.5,0.08),\ (0,0.30),
 \qquad p=(0.45,0.35,0.20).
\]
Both this law and the independent coupling have
\(\E F=0.625\) and \(\E K=0.097\). Nevertheless, the independent cost is
\(0.035750\), whereas the paired cost is \(0.073375\): the preferred
choice changes from passive to immediate execution. The loss from using
the independently selected action under the paired law is \(0.013375\)
per unit parent quantity. Under a convex mixture with paired-law weight
\(\theta\), passive cost is
\(0.035750+0.037625\theta\), crossing the market cost at approximately
\(\theta=0.64452\). These are normalized synthetic costs, not market bps.

\begin{table}[htbp]
\centering\small
\caption{Same marginal fills and completion prices; different joint laws.
The marginal-only worst case is an extremum of
\eqref{eq:completiontransport}. All values are synthetic.}
\begin{tabular}{lrr}
\toprule
Law & Passive expected cost & Selected action\\
\midrule
\inputtable{tables/realism_rows.tex}
\bottomrule
\end{tabular}
\end{table}

The calculation in \texttt{scripts/realism\_stress.py} compares exact
rational expectations with the marginal-coupling linear program and
checks the routing switch on a mixture grid. It is a dependence stress,
not an identification of endogenous impact. If \(K\) also changes with
submitted size or timing, that response requires its own controlled law.

\subsection{A falsifiable evaluation contract}

Section~\ref{sec:contract} turns these requirements into a preregistered
evaluation contract.
